% !TeX program = lualatex
% !TeX TXS-program:compile = txs:///lualatex/[--shell-escape]
% Copyright (c) 2026 A. Nader. Part of cacography, under the LaTeX Project Public
% License 1.3c or later (https://www.latex-project.org/lppl.txt).
% Bug reports: https://github.com/aynader/cacography/issues
% Demo: linear algebra in pencil on graph paper, for comparing handwriting models.
\documentclass{article}
\usepackage[paper=graph, ink=pencil, sloppiness=0.55, seed=2024]{cacography}
\usepackage{amsmath}

\begin{document}

\section{Eigenvalues, a quick recap}
A square matrix $A$ has an eigenvector $v$ when multiplying by it only stretches the
vector, so that $Av = \lambda v$ for some number $\lambda$. Moving everything to one side
gives a homogeneous system, and it has a nonzero solution only when
\begin{equation}
  \det(A - \lambda I) = 0 .
\end{equation}
This is the characteristic polynomial. For a two by two matrix it is a quadratic, so
there are at most two eigenvalues, and they can be complex even when every entry of the
matrix is real.

\subsection{Things to remember for the exam}
\begin{itemize}
  \item The trace is the sum of the eigenvalues.
  \item The determinant is their product.
  \item A symmetric matrix has real eigenvalues and orthogonal eigenvectors.
\end{itemize}
If the eigenvectors span the whole space we can write $A = PDP^{-1}$, where the columns
of $P$ are the eigenvectors and $D$ is diagonal. Powers then become easy, because
\begin{equation}
  A^k = P D^k P^{-1}
\end{equation}
and raising a diagonal matrix to a power only raises each entry. Ask the tutor on Friday
why this fails for a shear, which has only one direction that stays fixed.

\end{document}
